\clearpage
\phantomsection
\section*{Glossary}
\addcontentsline{toc}{section}{Glossary}

\begin{description}
\raggedright
\sloppy
\setlength{\itemsep}{0.2pt}
\setlength{\parsep}{0pt}
\setlength{\parskip}{0pt}
\item[Alfv\'{e}n Eigenmode (AE)] Discrete electromagnetic oscillation excited within shear-Alfv\'{e}n continuum spectral gaps created by periodic toroidal or helical variations in the confining magnetic field strength.

\item[Analytic Signal] Complex-valued time-series representation $z(t) = x(t) + \mathrm{i}\,\mathcal{H}[x(t)]$ formed from a real fluctuation signal $x(t)$ and its Hilbert transform $\mathcal{H}[x(t)]$, from which the instantaneous amplitude envelope $A(t)$ ($\mathrm{V}$) and unwrapped phase $\phi(t)$ ($\mathrm{rad}$) are extracted.

\item[Auto-bicoherence ($b^2$)] Dimensionless normalized magnitude-squared bispectrum bounded on $[0, 1]$ that quantifies the fraction of spectral power at a coupled frequency $f_3 = f_1 \pm f_2$ ($\mathrm{Hz}$) attributable to quadratic phase locking rather than independent random fluctuations.

\item[Bessel Filter] Analog or digital linear filter designed for maximally flat group delay in the passband, preserving pulse waveforms and instantaneous phase trajectories without ringing distortion across the $40\text{--}80$ and $80\text{--}120~\mathrm{kHz}$ mode bands.

\item[Beta-induced Alfv\'{e}n Eigenmode (BAE)] Low-frequency compressional shear-Alfv\'{e}n mode located below the geodesic acoustic mode (GAM) continuum gap whose frequency scales with the plasma ion-acoustic velocity ($\propto \sqrt{\tfrac{7}{4}T_i + T_e}$, with $T_e, T_i$ in $\mathrm{K}$ or $\mathrm{keV}$) and finite thermal pressure $\beta$.

\item[Bispectrum ($B$)] Third-order spectral moment $B(f_1, f_2) = \langle X(f_1) X(f_2) X^*(f_1 \pm f_2) \rangle$ ($\mathrm{V^3/Hz^3}$) probing coherent three-wave quadratic phase coupling across resonant sum or difference frequency triads $f_3 = f_1 \pm f_2$ ($\mathrm{Hz}$).

\item[Circular Concentration Coefficient ($r_{\mathrm{circ}}$)] Directional statistics metric bounded on $[0, 1]$ measuring angular clustering of spatial phase residuals across a magnetic probe array for candidate poloidal mode number $m$, where $r_{\mathrm{circ}} = 1$ indicates zero phase dispersion.

\item[Continuous Wavelet Transform (CWT)] Multiresolution time--frequency decomposition convolving a fluctuation signal against dilated copies of an analyzing wavelet (complex Morlet $\psi(t)$ in this thesis) to resolve rapid frequency chirps without fixed-window resolution trade-offs.

\item[Continuum Gap] Spectral frequency window in the continuous shear-Alfv\'{e}n spectrum opened by periodic toroidal ($R_0$) or helical ($L=1, M=4$) magnetic field variations, allowing discrete eigenmodes (such as AEs, BAEs, and TAEs) to persist with reduced continuum damping.

\item[Diamagnetic Loop] Inductive flux coil encircling the plasma cross-section to measure the excluded toroidal magnetic flux $\Delta\Phi_{\mathrm{dia}}$ ($\mathrm{Wb}$) proportional to perpendicular kinetic pressure, yielding the macroscopic diamagnetic stored energy $W_p$ ($\mathrm{J}$).

\item[Effective Sample Size ($N_{\mathrm{eff}}$)] Autocorrelation-adjusted count of statistically independent observations derived via Bartlett's variance inflation factor, preventing inflated significance in smoothed or serially correlated plasma profiles.

\item[Electron Cyclotron Emission (ECE)] Microwave radiation emitted by gyrating electrons at harmonics of the electron cyclotron frequency, measured at the second extraordinary (X2) harmonic ($57.5\text{--}72.5~\mathrm{GHz}$) to localize core electron-temperature fluctuations ($T_e$, $\mathrm{K}$, $\mathrm{keV}$, or uncalibrated relative intensity).

\item[Electron Cyclotron Heating (ECH)] Auxiliary radio-frequency plasma heating scheme using a $70~\mathrm{GHz}$ gyrotron resonant at the fundamental or second-harmonic electron cyclotron layer to heat electrons selectively.

\item[Energetic Particle Mode (EPM)] Kinetic magnetohydrodynamic instability driven directly by wave--particle resonances with super-Alfv\'{e}nic fast ions when kinetic drive exceeds continuum damping, often exhibiting rapid frequency chirping tied to fast-ion phase-space evolution.

\item[False Discovery Rate (FDR-BH)] Benjamini--Hochberg multiple-hypothesis adjustment controlling the expected fraction of false-positive correlations or triad detections across simultaneous statistical tests.

\item[Fast Fourier Transform] Computational algorithm for evaluating the discrete Fourier transform of $N_{\mathrm{fft}}$-sample time-series segments to obtain complex frequency spectra $X(f)$ ($\mathrm{V/Hz}$) for STFT spectrograms, Welch coherence, and bispectral estimation.

\item[Geodesic Acoustic Mode (GAM)] Axisymmetric electrostatic and compressible fluid oscillation driven by geodesic magnetic curvature whose frequency scales with the plasma ion-acoustic speed ($\propto \sqrt{T_e}$) and defines the lower accumulation boundary for BAE gaps.

\item[Global Alfv\'{e}n Eigenmode (GAE)] Discrete shear-Alfv\'{e}n oscillation localized near an extremum of the Alfv\'{e}n continuum in low-shear helical or tokamak configurations without requiring toroidicity-induced gap coupling.

\item[Heliotron] Non-axisymmetric toroidal stellarator-type magnetic confinement device using continuous helical coils alongside toroidal and vertical field coils to generate rotational transform ($\iota$) without requiring a large net ohmic plasma current ($I_p$, $\mathrm{A}$, typically in the $\mathrm{kA}$ range).

\item[High-Confinement Mode (H-mode)] Improved toroidal plasma confinement regime featuring reduced edge radial transport, steep pedestal pressure gradients, and enhanced stored energy $W_p$ ($\mathrm{J}$) above a heating power threshold.

\item[Hilbert Transform ($\mathcal{H}$)] Linear integral operator shifting all positive-frequency Fourier components by $-\pi/2~\mathrm{rad}$ to construct the quadrature component $\mathcal{H}[x(t)]$ of the complex analytic signal $z(t)$ for instantaneous envelope $A(t)$ and phase $\phi(t)$ extraction.

\item[Last Closed Flux Surface (LCFS)] Outermost nested toroidal magnetic surface bounding the confined plasma column ($\rho = r/a = 1.0$, $a = 0.170~\mathrm{m}$ in Heliotron~J) before field lines intersect the vacuum vessel or divertor.

\item[Lawson Criterion] Ignition condition for self-sustained fusion requiring the triple product of bulk density, energy confinement time, and thermal temperature to satisfy $n \tau_E T \ge 3 \times 10^{21}\,\mathrm{m^{-3}\cdot s\cdot keV}$.

\item[Magnetohydrodynamics (MHD)] Single-fluid macroscopic plasma framework combining Maxwell's equations with continuum fluid mechanics to describe magnetic equilibrium force balance and collective electromagnetic plasma waves.

\item[Magnitude-Squared Coherence ($\gamma^2$)] Dimensionless normalized cross-spectral correlation bounded on $[0, 1]$ measuring linear phase consistency between two probe signals (or between ECE and Mirnov channels) as a function of frequency $f$ ($\mathrm{Hz}$).

\item[Median Absolute Deviation ($\hat{\sigma}_{\mathrm{MAD}}$)] Outlier-resistant dispersion estimator given by $1.4826 \times \mathrm{median}(|x_i - \tilde{x}|)$ ($\mathrm{V}$, where $\tilde{x} = \mathrm{median}(x)$), providing a calibrated Gaussian-equivalent noise-floor estimate for adaptive burst detection.

\item[Mirnov Coil] Inductive pickup probe outside the last closed flux surface (LCFS, area-turns $N A$, $\mathrm{m^2}$) measuring $\mathrm{d}\tilde{B}_\theta/\mathrm{d}t$ ($\dot{B}$) fluctuations ($\mathrm{T}/\mathrm{s}$, amplified in $\mathrm{V}$) to resolve spectra and mode numbers ($m, n$).

\item[Neutral Beam Injection (NBI)] Auxiliary heating injecting $\sim 30~\mathrm{keV}$ neutral hydrogen beams that ionize to sustain super-Alfv\'{e}nic fast ions capable of destabilizing Alfv\'{e}nic modes.

\item[Non-Uniform Discrete Fourier Transform (NUDFT)] Spatial spectral transform evaluated over irregularly spaced coordinates (such as poloidal coil angles $\theta_j$, $\mathrm{rad}$, for PMP1--PMP14) to extract mode numbers $m$ without interpolation errors.

\item[Normalized Minor Radius ($\rho$)] Dimensionless flux-surface radial coordinate defined as $\rho = r/a$ relative to average minor radius $a$ ($\mathrm{m}$), spanning magnetic axis ($\rho = 0$) to the LCFS ($\rho = 1$).

\item[Optical Thickness ($\tau_{\mathrm{opt}}$)] Line-integrated absorption metric along an ECE chord; optically thick layers ($\tau_{\mathrm{opt}} \gg 1$) emit blackbody radiation proportional to electron temperature (uncalibrated intensity $I_{\mathrm{ECE}} \propto T_e$).

\item[Phase-Randomized Surrogate] Null-hypothesis series retaining the exact experimental Fourier amplitudes while drawing uniform random phases on $[0, 2\pi)$, destroying three-wave phase locking while preserving power spectra and $O(1/K)$ bias.

\item[Rotational Transform ($\iota$)] Magnetic field-line winding angle ($\mathrm{rad}$ per toroidal transit, normalized as $\iota/2\pi = 1/q$) quantifying helical twist around the torus and governing shear-Alfv\'{e}n continuum gap frequencies.

\item[Safety Factor ($q$)] Dimensionless inverse measure of rotational transform defined as $q = 2\pi/\iota$, representing the average number of toroidal field-line transits per single poloidal transit around the magnetic axis.

\item[Savitzky--Golay Filter] Polynomial local least-squares smoothing and differentiation filter ($525$-sample, $0.525~\mathrm{ms}$ window in this thesis) applied to unwrapped analytic phases $\phi(t)$ to extract instantaneous mode frequencies $f(t)$ while suppressing high-frequency phase-derivative noise.

\item[Sawtooth Oscillation] Quasi-periodic core relaxation instability producing rapid repetitive collapses and gradual recoveries in central electron temperature $T_e$ and density profiles via magnetic reconnection near low-order rational surfaces.

\item[Short-Time Fourier Transform (STFT)] Spectrogram method evaluating windowed discrete Fourier transforms over sliding time segments of duration $T$ ($\mathrm{s}$) with uniform frequency resolution $\Delta f = 1/T$ ($\mathrm{Hz}$).

\item[Three-Wave Coupling] Nonlinear quadratic interaction among wave modes satisfying frequency matching $f_3 = f_1 \pm f_2$ ($\mathrm{Hz}$) and coherent phase locking, enabling cross-band energy exchange and harmonic steepening quantified by bispectral bicoherence.

\item[Toroidal Alfv\'{e}n Eigenmode (TAE)] Shear-Alfv\'{e}n discrete mode residing within the spectral gap formed by toroidal coupling of adjacent poloidal harmonics ($m$ and $m+1$) centered near $f_{\mathrm{TAE}} = (\iota/2\pi)v_A/(4\pi R_0)$ ($\mathrm{Hz}$).
\end{description}

\clearpage
\phantomsection
\setcounter{section}{0}
\renewcommand{\thesection}{\Alph{section}}
\renewcommand{\theHsection}{appendix.\thesection} % <--- Unique internal link ID
\renewcommand{\theHsubsection}{appendix.\thesection.\arabic{subsection}}
\myappendix{Python Pipeline and GitHub Repository}
\label{appendix:github}

All signal-processing scripts, statistical auditing routines, and figure-generation tools developed in the framework of this thesis are publicly hosted in an open-source GitHub repository:\par\smallskip\noindent
{\footnotesize\url{https://github.com/Kevin-Rishor/Analysis-of-MHD-oscillations-at-Heliotron-J-with-Python-}}

The repository preserves the reproducible modular workflow built upon the \texttt{jpack} data-access layer and standard scientific Python libraries.

\subsection{General Pipeline Architecture}
\begin{sloppypar}
The codebase is organized under \texttt{analysis/} into an objective-driven modular hierarchy that transforms raw Heliotron~J \texttt{.edf} records into validated spectral, spatial, correlation, and higher-order spectral diagnostics without manual shot-by-shot tuning. Shared signal conditioning, coordinate mappings, and statistical utilities reside in \texttt{analysis/\allowbreak common/} (\texttt{mhd\_\allowbreak common.py} and \texttt{abicoh2\_\allowbreak corrected.py}) and \texttt{analysis/\allowbreak utils/} (\texttt{heliotron\_\allowbreak j\_\allowbreak ece.py}). Dedicated execution directories implement each specific objective: \texttt{analysis/\allowbreak obj1\_\allowbreak mode\_\allowbreak identification/} executes time--frequency spectrograms, poloidal/toroidal mode extraction, and multi-channel ECE radial localization; \texttt{analysis/\allowbreak obj2\_\allowbreak dynamics\_\allowbreak and\_\allowbreak structure/} executes Hilbert envelope demodulation, adaptive MAD burst gating, autocorrelation-adjusted correlation analyses, and $T_e$-proxy validation; \texttt{analysis/\allowbreak obj3\_\allowbreak nonlinear\_\allowbreak coupling/} executes auto-bicoherence mapping, split-sample $1/K$ asymptotic regression, and Fourier phase-randomized surrogate testing; and \texttt{analysis/\allowbreak obj\_\allowbreak ech/} executes automated ECH switch-off transition synthesis across discharge windows.
\end{sloppypar}

\subsection{Core Modules and Signal Processing Routines}
\begin{sloppypar}
Within this architecture, the main orchestration scripts execute end-to-end diagnostics for each research objective. Specifically, \texttt{mhd\_\allowbreak analysis\_\allowbreak obj1.py}, \texttt{mhd\_\allowbreak primary\_\allowbreak mode\_\allowbreak id.py}, and \texttt{analyze\_\allowbreak ece\_\allowbreak all\_\allowbreak probes\_\allowbreak radial.py} compute Hann-windowed STFT and complex Morlet CWT spectrograms paired with $3\sigma$ noise-floor gating and $16$-channel ECE--Mirnov cross-coherence profiles. Envelope dynamics and dual-harmonic mode identification are implemented in \texttt{mhd\_\allowbreak analysis\_\allowbreak obj2.py}, \texttt{mhd\_\allowbreak two\_\allowbreak mode\_\allowbreak combination\_\allowbreak search.py}, and \texttt{verify\_\allowbreak ece\_\allowbreak lpf\_\allowbreak temperature.py}, combining Bessel bandpass filtering, Savitzky--Golay phase differentiation, $N_{\mathrm{eff}}$-adjusted and partial correlations, and $30~\mathrm{Hz}$ low-pass ECE electron-temperature proxy verification against diamagnetic measurements. Nonlinear three-wave coupling routines reside in \texttt{abicoh2\_\allowbreak corrected.py}, \texttt{mhd\_\allowbreak analysis\_\allowbreak obj3.py}, and \texttt{mhd\_\allowbreak plot\_\allowbreak asymptotic\_\allowbreak bias\_\allowbreak scaling.py}, providing the index-corrected auto-bicoherence estimator and $500$-surrogate asymptotic convergence validation, while \texttt{run\_\allowbreak ech\_\allowbreak synthesis.py} and \texttt{compute\_\allowbreak bicoherence\_\allowbreak maps\_\allowbreak and\_\allowbreak controls.py} compile transition-window tables and cross-discharge reproducibility checks.
\end{sloppypar}

\myappendix{Adapted Poloidal Mode Number Extraction Script}
\label{appendix:poloidalm}
Listing~\ref{lst:poloidal} presents the author's adapted Python poloidal mode-extraction script (\texttt{Code/\allowbreak hjmag\_\allowbreak poloidal.py}), built upon the \texttt{turnelib} data-ingestion layer from the \texttt{jpack} toolkit (Annex~\ref{annex:library}) originally created by Prof.\ Shigeru Inagaki, Ph.D. Compared with the distributed baseline utilities, this script implements hardware-calibrated voltage-to-$\mathrm{T}/\mathrm{s}$ conversion ($N A = 4.5 \times 10^{-3}~\mathrm{m^2}$, $G = 2.0$), software polarity inversion on poloidal coils PMP1--PMP4 to correct reversed cabling, a non-uniform discrete Fourier transform (NUDFT) over the irregular angular positions of PMP1--PMP14, an exhaustive two-mode coexisting grid search across all $\binom{41}{2} = 820$ candidate harmonic pairs $(m_1, m_2) \in [-20, +20]^2$, and directional phase validation via the Mardia--Jupp circular concentration coefficient ($r_{\mathrm{circ}} \ge 0.70$) and mean absolute phase residual ($\le 45^\circ$).

\lstinputlisting[
    style=pythonstyle,
    basicstyle=\ttfamily\scriptsize,
    lineskip=-0.6pt,
    aboveskip=4pt,
    belowskip=2pt,
    caption={Adapted Heliotron~J poloidal mode number extraction script implementing polarity inversion, calibration, NUDFT, dual-harmonic grid search, and circular concentration validation.},
    label={lst:poloidal}
]{Code/hjmag_poloidal.py}
